%==================================================================
% BAB 5 — KESIMPULAN DAN SARAN
%
% Kesimpulan disusun BERNOMOR SEJAJAR dengan rumusan masalah:
% butir 1 menjawab RQ1, butir 2 menjawab RQ2, butir 3 menjawab RQ3.
% Jangan menambahkan butir kesimpulan yang tidak menjawab rumusan
% masalah mana pun, dan jangan meninggalkan rumusan masalah tanpa
% kesimpulan.
%==================================================================

\chapter{KESIMPULAN DAN SARAN}

\section{Kesimpulan}
Berdasarkan hasil penelitian dan pembahasan, diperoleh kesimpulan berikut.

\begin{daftarnomor}
    \item Kualitas layanan SIA berada pada kategori cukup hingga baik, dengan
          dimensi \textit{usability} memperoleh nilai kinerja terendah sebesar
          3,12 dan \textit{information quality} tertinggi sebesar 3,68.
          \emph{(menjawab \cref{rq:kualitas})}
    \item Seluruh dimensi memiliki kesenjangan negatif antara kepentingan dan
          kinerja. Kesenjangan terbesar terjadi pada dimensi \textit{usability}
          sebesar $-1{,}29$, dengan butir kemudahan navigasi sebagai
          penyumbang terbesar. \emph{(menjawab \cref{rq:kesenjangan})}
    \item Terdapat lima butir yang berada pada kuadran prioritas utama dan
          seluruhnya berkaitan dengan navigasi serta kejelasan alur, sehingga
          perbaikan disarankan diarahkan pada arsitektur informasi antarmuka.
          \emph{(menjawab \cref{rq:prioritas})}
\end{daftarnomor}

\section{Keterbatasan Penelitian}
% Subbab ini opsional pada sebagian program studi, namun sangat
% disarankan. Menyatakan keterbatasan secara jujur memperkuat
% kredibilitas, dan mendahului pertanyaan penguji.

\begin{daftarnomor}
    \item Penelitian dilakukan pada satu perguruan tinggi, sehingga hasilnya
          belum dapat digeneralisasi ke sistem informasi akademik lain.
    \item Pengukuran bersifat persepsional dan belum dilengkapi pengujian
          teknis terhadap antarmuka maupun kinerja sistem.
\end{daftarnomor}

\section{Saran}

\begin{daftarnomor}
    \item Bagi pengelola sistem, perbaikan disarankan dimulai dari kelima butir
          pada kuadran prioritas utama, khususnya penataan ulang struktur menu
          dan penyederhanaan alur pengisian rencana studi.
    \item Bagi penelitian selanjutnya, hasil pengukuran persepsional ini dapat
          dilengkapi dengan pengujian ketergunaan berbasis tugas untuk
          memperoleh bukti perilaku, bukan hanya persepsi.
    \item Penelitian serupa dapat dilakukan pada beberapa perguruan tinggi
          sekaligus agar diperoleh gambaran yang lebih dapat digeneralisasi.
\end{daftarnomor}
